\chapter{Área de Estudio y Datos}
\label{ch:datos}

Este capítulo describe el área de estudio, la unidad espacial de análisis y los datos que
alimentan la metodología. Un tema recurrente es la integridad de los datos: varios de los
resultados de esta tesis dependen de correcciones a la geografía base y a indicadores
específicos que se descubrieron durante una reconstrucción completa de reproducibilidad, y
se documentan aquí para que las cifras de todo el trabajo sean rastreables.

\section{La Zona Metropolitana de Guadalajara}
\label{sec:datos-zmg}

El área de estudio principal es la Zona Metropolitana de Guadalajara (\zmg), la segunda
más grande de México. Siguiendo la delimitación oficial, comprende diez municipios del
estado de Jalisco, que abarcan un núcleo central denso (Guadalajara y Zapopan) y una
periferia en rápida urbanización (\cref{tab:municipios}). Su movilidad se caracteriza por
una dependencia del automóvil alta y creciente, un extenso sistema de autobuses
concesionado y de operación informal y, más recientemente, una pequeña capa prémium de
autobús de tránsito rápido (Mi Macro) y tren ligero (Mi Tren). Esta capa prémium sigue el
patrón latinoamericano más amplio en que el autobús de tránsito rápido se ha expandido con
rapidez como complemento costo-efectivo del ferrocarril
\parencite{hidalgo2013brt, cervero2013brt}; injertarlo sobre una red base mayormente
concesionada es justamente el contexto que motiva un enfoque de planeación basado en la
demanda en lugar de uno que sigue a la oferta.

\begin{table}[htbp]
\centering
\caption{Los diez municipios de la \zmg y su número de \ageb{}s.}
\label{tab:municipios}
\begin{tabular}{lS[table-format=3.0]}
\toprule
Municipio & {\ageb{}s} \\
\midrule
Zapopan                       & 481 \\
Guadalajara                   & 404 \\
Tlajomulco de Zúñiga          & 341 \\
San Pedro Tlaquepaque         & 213 \\
Tonalá                        & 207 \\
El Salto                      & 115 \\
Zapotlanejo                   & 41  \\
Ixtlahuacán de los Membrillos & 34  \\
Acatlán de Juárez             & 29  \\
Juanacatlán                   & 16  \\
\midrule
\textbf{Total}                & \textbf{1881} \\
\bottomrule
\end{tabular}
\end{table}

\section{El \ageb como unidad de análisis}
\label{sec:datos-ageb}

Todos los indicadores se agregan al \ageb (\emph{área geoestadística básica}), la unidad
de enumeración censal definida por el INEGI. El \ageb es lo bastante fino para resolver la
variación intraurbana en demanda y acceso, y a la vez lo bastante grueso para que los
atributos censales se reporten a su nivel. El universo de análisis se restringe a los
\ageb{}s urbanos de los diez municipios de la \zmg: se excluyen las celdas cuyo
\code{CVE\_AGEB} lleva un sufijo alfabético (\ageb{}s rurales), así como todos los
municipios fuera de la delimitación metropolitana. Tras este filtro el universo es de
\num{1881} \ageb{}s (\cref{tab:municipios}) ---y no los \num{2068} usados en versiones
previas de este trabajo, discrepancia que se explica en \cref{sec:datos-integridad}.

\section{Fuentes de datos}
\label{sec:datos-fuentes}

La metodología separa deliberadamente los datos de \emph{Nivel~1} ---censales, económicos
y de red, disponibles para cualquier ciudad mexicana--- de los datos de oferta de
transporte, usados únicamente para medir la brecha de cobertura y auditar las rutas
existentes. La \cref{tab:fuentes-datos} lista las fuentes principales.

\begin{table}[htbp]
\centering
\caption{Principales fuentes de datos y su papel en la metodología.}
\label{tab:fuentes-datos}
\begin{tabularx}{\linewidth}{llX}
\toprule
Fuente & Nivel & Papel \\
\midrule
INEGI CPV2020    & \ageb / manzana & Población, proporción joven, propiedad vehicular \\
DENUE            & establecimiento & Proxies de empleo y atracción de POI \\
OpenStreetMap    & red             & Intersecciones, densidad vial, grafo vial \\
GTFS             & parada / ruta   & Accesibilidad existente y auditoría \\
EOD 2022         & zona de encuesta & Calibración del modelo gravitacional \\
CONAPO / CONEVAL & \ageb           & Marginación y rezago social (equidad) \\
\bottomrule
\end{tabularx}
\end{table}

\begin{description}[leftmargin=1.4em, style=nextline]
  \item[Censo y sociodemografía (INEGI CPV2020).] El Censo de Población y Vivienda 2020
    \parencite{inegi_cpv2020} aporta la población, la proporción de población joven usada
    para amplificar la generación de viajes y la propiedad vehicular de los hogares
    (\code{VPH\_AUTOM} sobre \code{VIVPAR\_HAB}) usada para derivar la propensión al
    transporte.
  \item[Actividad económica (DENUE).] El Directorio Estadístico Nacional de Unidades
    Económicas \parencite{inegi_denue} provee la ubicación de establecimientos, agregada a
    empleo por \ageb y a densidades de puntos de interés y comercio que impulsan la
    atracción de viajes y las oportunidades de accesibilidad.
  \item[Red vial (OpenStreetMap).] El grafo vial y los indicadores de Nodo derivados
    (densidad de intersecciones y de calles) provienen de \textsc{osm}
    \parencite{openstreetmap}, extraídos para los diez municipios con \textsc{osmnx}
    \parencite{boeing2017osmnx}.
  \item[Oferta de transporte (GTFS).] La fuente \textsc{gtfs} de SITEUR define la red
    existente usada para la superficie de accesibilidad (\cref{sec:w3}) y la auditoría de
    rutas (\cref{sec:w7}).
  \item[Encuesta Origen--Destino (EOD 2022).] La encuesta metropolitana de viajes 2022
    \parencite{imeplan_eod2022} provee flujos zonales observados para calibrar el
    decaimiento con la distancia del modelo gravitacional (\cref{sec:w2}).
  \item[Índices de equidad (CONAPO, CONEVAL).] El índice de marginación
    \parencite{conapo_marginacion2020} y el índice de rezago social
    \parencite{coneval_rezago2020} conforman el término de equidad (\cref{eq:equity}).
\end{description}

\section{Convenciones espaciales e integridad de datos}
\label{sec:datos-integridad}

Todo el cómputo usa el sistema de referencia proyectado canónico EPSG:6372 (cónico
equidistante para México); las coordenadas geográficas (EPSG:4326) se usan solo para la
ingesta. Todas las columnas de geometría están indexadas espacialmente.

Una reconstrucción completa de reproducibilidad, desde cero, sacó a la luz tres defectos de
integridad silenciosamente presentes en versiones previas, todos ya corregidos. Primero, un
error de dedo en el código de municipio del filtro de la geografía base había eliminado un
municipio entero en cada reconstrucción; corregirlo, junto con excluir correctamente los
\ageb{}s de sufijo alfabético, da el universo de \num{1881} \ageb{}s usado en todo el
trabajo (las cifras anteriores, computadas contra una tabla de \num{2068} filas nunca
reproducible, deben leerse como superadas). Segundo, el ráster de elevación digital se había
cargado bajo una referencia espacial equivocada, produciendo valores de pendiente
geométricamente sin sentido; el modelo heredado afectado se retiró en lugar de re-ejecutarse.
Tercero, el índice de marginación entraba al término de equidad con el signo invertido y con
un relleno de ceros fuera de rango para los \ageb{}s no emparejados, lo que en conjunto había
permitido que el bono de equidad premiara a las zonas \emph{menos} marginadas; corregir la
dirección y usar imputación por mediana restaura un término de equidad monótono y con el
signo correcto. Estas correcciones afectan de manera sustantiva varias cifras principales y
son la razón por la que los resultados de esta tesis superan cualquier cifra anterior para
las mismas cantidades.
